Глава~3 закончилась спуском с постоянным шагом $1/L$, для которого нужно знать константу $L$. Здесь шаг подбирается на каждой итерации по \hyperref[note:5330B23D-95DD-4D5D-B4A3-63EBADB1DB7D]{условиям Армихо--Вульфа}, а $d_k$ может быть любым направлением спуска. Сходимость держится на двух оценках: \hyperref[note:93FC9CC4-79FA-4CEA-9B10-9DBFB2C07850]{оценке на шаг} и \hyperref[note:E4AE9B25-BDB9-4A54-8CCB-09809F20B79D]{косинусе угла} между антиградиентом и $d_k$, отделённом от нуля; итог --- $\min_k \|\nabla f(x_k)\| = O(1/\sqrt{n})$. В конце главы появляется первое направление, отличное от антиградиента, --- \hyperref[note:C0ED5501-CFC6-42BA-A53A-1B346CBA545D]{метод Ньютона}.
